% Fuente única para el anexo incluido en la memoria y el fichero ODS independiente.
% Marcar exactamente una columna por fila con \textbf{X}.
% Categorías exigidas por la plantilla ETSINF: alto, medio, bajo y no procede.
\begin{center}
\resizebox{\textwidth}{!}{%
\begin{tabular}{|l|c|c|c|c|}\hline
\textbf{Objetivos de Desarrollo Sostenible} & \textbf{Alto} & \textbf{Medio} &
\textbf{Bajo} & \textbf{No} \\
& & & & \textbf{procede} \\ \hline
ODS 1. \textbf{Fin de la pobreza.} & & & & \textbf{X} \\ \hline
ODS 2. \textbf{Hambre cero.} & & & & \textbf{X} \\ \hline
ODS 3. \textbf{Salud y bienestar.} & & & & \textbf{X} \\ \hline
ODS 4. \textbf{Educación de calidad.} & & & \textbf{X} & \\ \hline
ODS 5. \textbf{Igualdad de género.} & & & & \textbf{X} \\ \hline
ODS 6. \textbf{Agua limpia y saneamiento.} & & & & \textbf{X} \\ \hline
ODS 7. \textbf{Energía asequible y no contaminante.} & & & & \textbf{X} \\ \hline
ODS 8. \textbf{Trabajo decente y crecimiento económico.} & & & & \textbf{X} \\ \hline
ODS 9. \textbf{Industria, innovación e infraestructuras.} & \textbf{X} & & & \\ \hline
ODS 10. \textbf{Reducción de las desigualdades.} & & & \textbf{X} & \\ \hline
ODS 11. \textbf{Ciudades y comunidades sostenibles.} & & \textbf{X} & & \\ \hline
ODS 12. \textbf{Producción y consumo responsables.} & & \textbf{X} & & \\ \hline
ODS 13. \textbf{Acción por el clima.} & & & \textbf{X} & \\ \hline
ODS 14. \textbf{Vida submarina.} & & & & \textbf{X} \\ \hline
ODS 15. \textbf{Vida de ecosistemas terrestres.} & & & & \textbf{X} \\ \hline
ODS 16. \textbf{Paz, justicia e instituciones sólidas.} & & & & \textbf{X} \\ \hline
ODS 17. \textbf{Alianzas para lograr los objetivos.} & & & & \textbf{X} \\ \hline
\end{tabular}%
}
\end{center}

\newpage

\section*{Reflexión}

La relación principal del TFG es con el ODS 9, Industria, innovación e
infraestructuras. El trabajo estudia la inteligencia artificial aplicada al patrimonio musical
digital y propone un procedimiento reproducible para decidir cuándo utilizarla. Incluye
degradación normalizada por escala de pentagrama, controles contra fugas de datos,
comparaciones emparejadas, revisión de fallos de notación y una guía profesional por
condición. Este componente metodológico y de innovación responsable justifica una relación alta,
sin desplegar un producto ni infraestructura de producción.

La relación con el ODS 11, Ciudades y comunidades sostenibles, es media. Su meta 11.4 contempla la
protección del patrimonio cultural \cite{un2030Agenda}, que incluye las partituras digitalizadas.
Una copia ampliada de consulta puede facilitar el acceso y reducir la manipulación del
original. La superresolución no conserva por sí misma un documento: puede alterar
símbolos, texto o trazos y generar una apariencia engañosa. Por ello, el uso recomendado mantiene
siempre la imagen original, identifica el derivado y exige cotejo humano cuando haya consecuencias
editoriales, interpretativas o de conservación. Se busca evitar confundir mejora visual y
restauración histórica auténtica.

El ODS 12, Producción y consumo responsables, tiene relación media. Se compararon
fidelidad y coste temporal: SwinIR fue entre 5,2 y 8,5 veces más lento que EDSR en la ejecución
retenida, sin ventaja general. Su uso indiscriminado habría consumido más recursos sin beneficio
consistente. El marco profesional propone EDSR como opción general y reserva SwinIR para las
condiciones donde su ganancia emparejada justifica el coste. También distingue el acceso a datos
de su reproducción: SMB se utiliza con acceso autorizado y no se redistribuye el corpus bruto;
la memoria incorpora comparaciones analíticas seleccionadas, atribuidas y con
transformaciones identificadas. Cada licencia y la base de publicación de cada imagen se consideran por
separado, sin presumir que el acceso autorice cualquier reproducción.

La relación con el ODS 13, Acción por el clima, es baja porque no se midió energía ni huella de
carbono con instrumentación validada. Se registraron hardware y duración y se limitaron modelos y
ejecuciones. Una comparación controlada, reutilizar código en local y Kaggle y acotar la adaptación
a un estudio secundario redujeron cálculo innecesario. Estas medidas ahorran recursos,
pero no permiten cuantificar un beneficio climático ni justifican una valoración mayor.

Los ODS 4 y 10 tienen relación baja e indirecta. Las copias de consulta podrían facilitar el
estudio de fuentes musicales a estudiantes, intérpretes o investigadores y mejorar el acceso a
material digitalizado de resolución insuficiente. El TFG no ha medido
aprendizaje, accesibilidad, cobertura geográfica ni reducción de desigualdades, y tampoco ha creado
un servicio público. La contribución es un escenario profesional futuro, no
impacto demostrado.

Los demás ODS no son aplicables. El
trabajo no aborda pobreza, alimentación, salud, género, agua, energía asequible, empleo, ecosistemas,
instituciones o alianzas mediante resultados directos. La innovación debe evaluarse junto con
sus riesgos: en patrimonio musical, un detalle falso puede ser
más perjudicial que una imagen borrosa porque parece fiable. La transparencia de las degradaciones,
los hashes, las licencias, los resultados negativos y los casos de no uso es
esencial para la responsabilidad social de la solución.
